\documentclass[a4paper,12pt]{article}

\usepackage[utf8]{inputenc}
\usepackage[T1]{fontenc}
\usepackage[brazil]{babel}
\usepackage{amsmath}
\usepackage{amssymb}
\usepackage{siunitx}
\usepackage{geometry}
\geometry{a4paper, left=3cm, right=2cm, top=3cm, bottom=2cm}
\usepackage{setspace}
\usepackage{indentfirst}
\usepackage{graphicx}
\usepackage{float} 
\usepackage{booktabs}
\usepackage{array}
\usepackage{tabularx}
\usepackage{enumitem}
% Macro para destacar trechos sem citação (autoavaliar depois): negrito + sublinhado
% Usa apenas comandos padrão LaTeX para máxima compatibilidade
\newcommand{\autoeval}[1]{\textbf{\underline{#1}}}
\usepackage{hyperref}
\hypersetup{
    colorlinks=true,
    linkcolor=black,
    filecolor=magenta,      
    urlcolor=blue,
    citecolor=black,
    pdfencoding=auto
}

\title{\textbf{PIBITI}\\ °\\Análise Comparativa de Motores e Hélices de uso em Aeronaves Agrícolas de Asa Fixa}
\author{Vinicius Andrade Trento}
\date{\today}
\onehalfspacing

\begin{document}

\maketitle
\newpage
\tableofcontents
\newpage


\section{Contextualização e Definição dos Requisitos de Engenharia}

O presente relatório constitui a execução da fase de pesquisa fundamental delineada no Plano de Trabalho PIBIT, com foco específico no desenvolvimento de soluções para o projeto NAPI - Aeronaves de Pequeno Porte. A missão central consiste em subsidiar a engenharia de sistemas para o AGRO-VANT, uma plataforma de asa fixa destinada a operações agrícolas que exigem robustez para cargas úteis elevadas e autonomia estendida.

A transição das configurações multirrotor para arquiteturas de asa fixa apresenta-se como uma evolução técnica necessária para maximizar a cobertura de área por voo e otimizar a eficiência energética. Essa mudança justifica-se pela limitação dos multirrotores em satisfazer requisitos de larga escala, dada a sua dependência de sustentação vetorial contínua. Contudo, à medida que a plataforma escala em massa máxima de decolagem (MTOW), o problema deixa de ser apenas de seleção de componentes COTS (\textit{Commercial Off-The-Shelf}) e passa a se inserir em um regime de engenharia aeronáutica certificada, com requisitos termodinâmicos, estruturais e regulatórios substancialmente distintos dos VANTs convencionais.

\subsection{Definição do Escopo: Aeronaves Agrícolas Leves Tripuladas (1{,}0 a 2{,}5 t)}

O presente trabalho concentra-se na faixa de Peso Máximo de Decolagem (MTOW) entre \SI{1000}{kg} e \SI{2500}{kg}, segmento que reúne as aeronaves agrícolas tripuladas de asa fixa mais utilizadas globalmente e que constitui referência direta para o dimensionamento do AGRO-VANT em sua configuração de máxima capacidade operacional. Esse recorte foi escolhido pelos seguintes motivos:

\begin{itemize}
    \item \textbf{Maturidade tecnológica e disponibilidade de dados:} A faixa abriga plataformas operacionais consolidadas há mais de cinco décadas (Piper PA-25 Pawnee, Cessna 188 AgWagon, Piper PA-36 Pawnee Brave, Embraer EMB-200/201/202 Ipanema), o que assegura uma base extensa de dados de campo, manuais técnicos e referências normativas para análise comparativa.
    \item \textbf{Convergência tecnológica em motor de pistão de grande cilindrada:} Diferentemente da faixa abaixo de \SI{600}{kg} (dominada por motores 2T COTS) ou da faixa acima de \SI{4000}{kg} (dominada por turboélices PT6A), o segmento 1--2{,}5 t apresenta soluções tecnologicamente convergentes em motores aeronáuticos de pistão de 6 a 8 cilindros horizontalmente opostos (boxer), com potências entre \SI{235}{HP} e \SI{400}{HP}.
    \item \textbf{Aderência ao mercado brasileiro:} O EMB-202 Ipanema (MTOW \SI{1800}{kg}) detém aproximadamente 80\% da frota de pulverização aérea brasileira, situando-se exatamente no centro da faixa estudada.
    \item \textbf{Quadro regulatório bem estabelecido:} Toda a faixa é regida pelos mesmos regulamentos (FAR 23 / RBAC 23 para aeronavegabilidade, FAR 137 / RBAC 137 para operação), permitindo um tratamento unificado dos requisitos de certificação.
\end{itemize}

A Tabela \ref{tab:plataformas_referencia} apresenta as principais plataformas de referência consideradas neste estudo.

\begin{table}[H]
\centering
\caption{Plataformas Agrícolas de Referência na Faixa de MTOW de 1 a 2{,}5 toneladas}
\label{tab:plataformas_referencia}
\footnotesize
\renewcommand{\arraystretch}{1.3}
\begin{tabularx}{\textwidth}{@{} 
    >{\hsize=1.0\hsize\raggedright\arraybackslash}X 
    c 
    c 
    >{\hsize=1.0\hsize\raggedright\arraybackslash}X 
    @{}}
\toprule
\textbf{Aeronave} & \textbf{MTOW (kg)} & \textbf{Hopper (L)} & \textbf{Motor (Potência)} \\ \midrule
Piper PA-25-235 Pawnee B & $\sim$1.315 & 568--680 & Lycoming O-540-B2B5 (235 HP) \\
\addlinespace
Cessna 188 AgWagon B & 1.497 & 757 & Continental IO-520-D (300 HP) \\
\addlinespace
Embraer EMB-200 Ipanema & $\sim$1.550 & 580 & Lycoming O-540-H2B5D (260 HP) \\
\addlinespace
Embraer EMB-201A Ipanema & $\sim$1.550 & 750 & Lycoming IO-540-K1D5 (300 HP) \\
\addlinespace
Embraer EMB-202 Ipanema & 1.800 & 950 & Lycoming IO-540-K1J5 (300/320 HP) \\
\addlinespace
Cessna A188B AgTruck / AgHusky & 1.996 & 1.060 & Continental IO-520-D / TSIO-520-T (300/310 HP) \\
\addlinespace
Piper PA-36 Pawnee Brave 300 & 1.996--2.177 & 850--1.080 & Lycoming IO-540-K1G5 (300 HP) \\
\addlinespace
Piper PA-36 Brave 375 / 400 & $\sim$2.177 & 1.080 & Lycoming IO-720-D1CD (375--400 HP) \\ \bottomrule
\multicolumn{4}{l}{\footnotesize Compilação do autor a partir de catálogos dos fabricantes e referências históricas \cite{embraer202, piper_pa25, cessna_188, piper_pa36}.}
\end{tabularx}
\end{table}

Observa-se que a faixa apresenta dois subgrupos identificáveis: as \textbf{aeronaves de 1{,}0--1{,}5 t}, predominantemente equipadas com motores Lycoming O-540 de \SI{235}{HP} a \SI{260}{HP} ou Continental IO-520 de \SI{300}{HP}, e as \textbf{aeronaves de 1{,}5--2{,}5 t}, equipadas com Lycoming IO-540 (\SI{300}{HP}) ou IO-720 (\SI{375}{HP} a \SI{400}{HP}).

\subsection{O Envelope Operacional Agrícola}
O envelope operacional agrícola impõe exigências específicas sobre a engenharia aeronáutica da plataforma:

\begin{itemize}
    \item \textbf{Ciclos de Carga e Dinâmica de Voo:} A operação caracteriza-se por ciclos de carga mecânica e térmica elevados. A aplicação envolve voos a baixa altitude sobre a cultura, exigindo resposta rápida do motor (\textit{throttle response}) para manobras de retorno e evasão de obstáculos imprevistos. Em motores de pistão de grande cilindrada (6 ou 8 cilindros formato boxer), essa resposta é favorecida pela inércia rotacional reduzida quando comparada a turboélices, característica que justifica a permanência da arquitetura de pistão nessa faixa.
    
    \item \textbf{Atmosfera Operacional:} As missões ocorrem em ambientes com alta concentração de material particulado em suspensão e vapores químicos provenientes dos defensivos. Essa condição exige índices de proteção IP (\textit{Ingress Protection}) elevados para os componentes eletrônicos, sistemas eficientes de filtragem de ar e selagem da cabine para casos em que é necessário pilotos (sistema característico do Ipanema e do AgWagon).
    
    \item \textbf{Densidade de Altitude:} Operações no interior do Brasil frequentemente ocorrem em dias quentes e regiões de planalto, resultando em redução da densidade do ar ($\rho$). Esse fator afeta a sustentação aerodinâmica e a eficiência volumétrica dos motores de combustão interna. Como regra prática estabelecida na literatura de desempenho aeronáutico, um motor naturalmente aspirado perde aproximadamente 3{,}5\% de potência a cada \SI{1000}{ft} de aumento na altitude-densidade \cite{cfisteph_density}, o que sustenta a recomendação de margem de potência instalada significativa em relação ao MTOW para garantir o desempenho na decolagem em pistas de planalto. A análise quantitativa do efeito da altitude sobre o desempenho será feita na seção \ref{sec:curvas_lycoming}, com base em dados do manual do fabricante.
    
    \item \textbf{Variação de Massa e Centro de Gravidade:} Durante a fase de pulverização, o peso da aeronave decresce substancialmente à medida que a carga útil é ejetada, alterando a posição do centro de gravidade (CG) e a carga alar (\textit{wing loading}). No Ipanema EMB-202, a ejeção de até \SI{950}{L} de produto representa proporção significativa do MTOW de \SI{1800}{kg} \cite{embraer202_wiki}; no Pawnee Brave 375, o hopper de aproximadamente \SI{1080}{L} (38 ft$^3$) representa proporção semelhante para o MTOW de \SI{2177}{kg} \cite{piper_pa36}.
\end{itemize}   

\section{Análise Comparativa de Motores na Faixa de Interesse}

A escolha do motor para aeronaves agrícolas tripuladas na faixa de 1 a 2{,}5 toneladas é fortemente condicionada pela carga útil pretendida, pela autonomia de missão e pelos requisitos de certificação. A Tabela \ref{tab:matriz_propulsao} sintetiza as três famílias dominantes de motores aeronáuticos de pistão certificados aplicáveis ao segmento.

\begin{table}[H]
\centering
\caption{Famílias de Motores Aeronáuticos de Pistão para Aeronaves Agrícolas de 1 a 2{,}5 toneladas}
\label{tab:matriz_propulsao}
\footnotesize
\renewcommand{\arraystretch}{1.3}
\begin{tabularx}{\textwidth}{@{} 
    >{\hsize=0.8\hsize\raggedright\arraybackslash}X 
    >{\hsize=0.9\hsize\raggedright\arraybackslash}X 
    c 
    >{\hsize=1.3\hsize\raggedright\arraybackslash}X 
    @{}}
\toprule
\textbf{Família} & \textbf{Configuração} & \textbf{Faixa de Potência} & \textbf{Aplicação Típica (MTOW)} \\ \midrule
Lycoming O-540 / IO-540 & 6 cil. boxer, 8{,}9 L (541{,}5 in$^3$) & 235--\SI{350}{HP} & PA-25 Pawnee (1{,}3 t), Ipanema (1{,}5--1{,}8 t), Brave 300 (2{,}0 t) \\
\addlinespace
Continental IO-520 / TSIO-520 & 6 cil. boxer, 8{,}5 L (520 in$^3$) & 285--\SI{310}{HP} & Cessna 188 AgWagon/AgTruck/AgHusky (1{,}5--2{,}0 t) \\
\addlinespace
Lycoming IO-720 & 8 cil. boxer, 11{,}8 L (722 in$^3$) & \SI{400}{HP} & PA-36 Brave 375/400 (2{,}2 t) \\ \bottomrule
\multicolumn{4}{l}{\footnotesize Dados consolidados a partir de \cite{lycomingmanual, io360manual, lycoming_io720}.}
\end{tabularx}
\end{table}

\section{Análise de Motores de Combustão Interna (ICE)}

Considerando o atual estado da arte do armazenamento eletroquímico de energia, os Motores de Combustão Interna (ICE - \textit{Internal Combustion Engines}) permanecem como a alternativa tecnicamente mais madura e operacionalmente viável para aeronaves agrícolas tripuladas na faixa de 1 a 2{,}5 toneladas. Essa predominância baseia-se na densidade energética específica dos combustíveis hidrocarbonetos, da ordem de \SI{12000}{Wh/kg} para o querosene de aviação \cite{battery_aviation_review, aerospace_america_batteries}, que supera as melhores baterias de íon-lítio comercialmente disponíveis (densidade típica de até \SI{330}{Wh/kg}) e mesmo baterias avançadas em estado sólido (até \SI{400}{Wh/kg}) \cite{battery_aviation_review} por mais de uma ordem de grandeza, \autoeval{tornando o vetor puramente elétrico inviável para missões de pulverização agrícola que demandam autonomia mínima de 2 horas com cargas úteis de \SI{500}{kg} a \SI{1000}{kg}}.


\subsection{Parâmetros Termodinâmicos e de Avaliação de Desempenho}
Para o dimensionamento da planta motriz, utiliza-se a Pressão Média Efetiva ao Freio (BMEP - \textit{Brake Mean Effective Pressure}) como métrica de otimização volumétrica e mecânica. A BMEP quantifica o trabalho efetuado pelo motor por ciclo, permitindo a comparação da eficiência volumétrica e mecânica independentemente da cilindrada \cite{heywood1988}. 

A autonomia (\textit{Endurance}) é avaliada pelo Consumo Específico de Combustível ao Freio (BSFC), que relaciona a vazão mássica de combustível consumida com a potência gerada, sendo a métrica principal para o dimensionamento do alcance de voo \cite{elsayed2016}. Adicionalmente, a Eficiência Térmica ao Freio (BTE) determina a capacidade do sistema em converter a energia química do combustível em trabalho mecânico útil. Para motores aeronáuticos de pistão refrigerados a ar, como a família Lycoming O-540, o BSFC em regime de cruzeiro situa-se na faixa de \SI{0.48}{lb/(hp \cdot h)} (aproximadamente \SI{292}{g/(kW \cdot h)}), valor próximo ao de turboélices modernos em regime otimizado \cite{lycomingmanual}.

\subsection{Arquitetura Boxer Horizontalmente Oposta: A Convergência da Faixa}

Os três principais motores aplicáveis à faixa de 1 a 2{,}5 t (Lycoming O-540, Continental IO-520 e Lycoming IO-720) compartilham a mesma arquitetura fundamental: pistões dispostos em configuração horizontalmente oposta (boxer), operando em ciclo Otto de 4 tempos, refrigeração a ar por aletas e acionamento direto da hélice (\textit{direct drive}). Essa convergência tecnológica é consequência de requisitos de engenharia compartilhados:

\begin{itemize}
    \item \textbf{Balanço Harmônico:} A disposição boxer permite o cancelamento dos momentos de inércia primários de primeira e segunda ordem, reduzindo vibrações transmitidas à estrutura e à aviônica embarcada. Esse aspecto é particularmente crítico em aplicações agrícolas, dada a fadiga estrutural acumulada por ciclos repetitivos de carga.
    \item \textbf{Refrigeração a Ar:} A ausência de circuito de líquido refrigerante simplifica a manutenção em campo, condizente com a operação descentralizada da aviação agrícola, em pistas rurais. Exige, contudo, defletores (\textit{engine baffling}) cuidadosamente dimensionados para garantir uniformidade térmica entre cilindros \cite{lycomingmanual, io360manual}.
    \item \textbf{Direct Drive:} O acoplamento direto entre virabrequim e hélice elimina a complexidade de uma caixa redutora, reduzindo peso, custo e modos de falha. A consequência é a operação do motor em faixa de RPM compatível com o regime aerodinâmico ótimo da hélice, tipicamente entre 2.400 e 2.700 RPM para os motores aeronáuticos das famílias O-540, IO-540 e IO-720 \cite{lycomingmanual, lycoming_io720}.
    \item \textbf{Eficiência Volumétrica e Consumo:} A separação mecânica das fases de admissão e exaustão via válvulas comandadas por tuchos hidráulicos maximiza o aproveitamento térmico, com BSFC em regime de cruzeiro entre 280 e \SI{350}{g/(kW \cdot h)} \cite{lycomingmanual, io360manual}.
    \item \textbf{Confiabilidade Certificada:} O Intervalo Entre Revisões Gerais (TBO) recomendado pelos fabricantes para a família é de 1.500 a 2.000 horas em serviço intermitente \cite{lycomingmanual}, valor adequado ao perfil de utilização típico de aeronaves agrícolas, considerando-se as características sazonais das safras e o uso intensivo durante a janela operacional \autoeval{(estimativa do autor: 200 a 400 h/ano em safra plena)}.
\end{itemize}

\subsection{Injeção de Combustível em Motores Aeronáuticos: Mecânica vs. Eletrônica}

Os motores aeronáuticos certificados da faixa de interesse adotam, em sua quase totalidade, sistemas de injeção mecânica contínua (IO --- \textit{Injected Opposed}). A injeção mecânica utiliza um servorregulador (tipicamente Bendix RSA-5AD1 ou RSA-10ED1) que mede a vazão de ar pela admissão e dosa proporcionalmente o combustível a bicos posicionados próximos às válvulas de admissão de cada cilindro \cite{io360manual}.

Comparada à carburação convencional, a injeção mecânica oferece vantagens importantes:

\begin{itemize}
    \item Imunidade ao \textit{carburetor icing} em condições de baixa temperatura e alta umidade.
    \item Distribuição mais uniforme da mistura entre cilindros, reduzindo gradientes de temperatura.
    \item Permite voo em atitudes negativas momentâneas (especialmente importante em manobras agrícolas evasivas).
    \item Resposta rápida do motor em transientes de potência (\textit{throttle response}), crítica em pulverização de baixa altitude.
\end{itemize}

A migração para sistemas eletrônicos de injeção e ignição (FADEC --- \textit{Full Authority Digital Engine Control}) é uma tendência emergente, exemplificada pela linha Lycoming iE2, que oferece controle eletrônico integral com sensores MAP (\textit{Manifold Absolute Pressure}) e IAT (\textit{Intake Air Temperature}). O FADEC ajusta automaticamente a razão estequiométrica da mistura combustível e ar, ângulo de ignição (pré-ignição) e mistura para densidade local do ar, eliminando a necessidade de ajuste manual de mistura pelo piloto. Para operações em densidade de altitude variável (planaltos brasileiros, com elevações de \SI{600}{m} a \SI{1100}{m}), essa correção automática representa ganho mensurável de eficiência operacional que provém de fundamentos de regimes de trabalho do motor: regime pobre, regime economico e regime rico (alto desempenho).



\section{Detalhamento das Aeronaves Agrícolas de Grande Porte}
\label{sec:classe4}

\subsection{A Família Lycoming O-540 e o Ipanema}

A aeronave agrícola de referência no Brasil é o Embraer/Neiva EMB-202 Ipanema, com mais de \num{1700} unidades entregues e participação de mercado próxima a 80\% da frota brasileira de pulverização aérea \cite{embraer202}. Suas especificações técnicas servem como \textit{benchmark} para qualquer plataforma de Classe IV nacional.

\begin{itemize}
    \item \textbf{MTOW:} \SI{1800}{kg}
    \item \textbf{Capacidade do hopper:} \SI{950}{L} (cerca de 53\% do MTOW)
    \item \textbf{Motor (versão EMB-202A a etanol):} Lycoming IO-540-K1J5, \SI{320}{HP} a \SI{2700}{RPM}, 6 cilindros boxer, deslocamento de \SI{8.9}{L} (\SI{541.5}{in^3})
    \item \textbf{Hélice:} Hartzell tripá de velocidade constante (passo variável)
    \item \textbf{TBO recomendado:} 1.500 horas
    \item \textbf{Capacidade de combustível utilizável:} \SI{264}{L}
\end{itemize}

A família Lycoming O-540 / IO-540, certificada pela FAA sob o \textit{Type Certificate} 1E4, é o motor de pistão dominante em aplicações agrícolas tripuladas dessa classe \cite{lycomingmanual}. Trata-se de um motor de 6 cilindros horizontalmente opostos (boxer), refrigerado a ar, com taxa de compressão entre 7{,}2:1 e 8{,}5:1. O peso seco situa-se em torno de \SI{438}{lb} (\SI{198.7}{kg}) e o BSFC em regime de cruzeiro é de aproximadamente \SI{0.48}{lb/(hp \cdot h)} (\SI{292}{g/(kW \cdot h)}) \cite{lycomingmanual}.

\subsubsection{O Caso do Etanol Aeronáutico}

Uma particularidade altamente relevante para o cenário brasileiro é a certificação do EMB-202A, a primeira aeronave de série do mundo a operar com motor a etanol \cite{embraer202}. A conversão é regulamentada pela ANAC por meio da Instrução Suplementar IS 137.201-001E, que exige a emissão de uma Autorização Especial de Voo (AEV). O etanol oferece vantagens econômicas significativas em relação à AvGas 100LL, além de eliminar o uso de chumbo tetraetila (metal pesado tóxico), mas demanda modificações no sistema de injeção e nos materiais da linha de combustível devido à sua agressividade química e menor densidade energética. O etanol apresenta cerca de 69\% da energia volumétrica do JP8 (querosene aeronáutico) \cite{nas_army_fuels}, o que implica em maior consumo volumétrico para entregar a mesma energia ao motor.

\subsection{Análise Quantitativa das Curvas Lycoming IO-360}
\label{sec:curvas_lycoming}

A documentação técnica dos motores Lycoming fornece dados quantitativos detalhados em duas curvas fundamentais que se aplicam, com adaptações de escala, a toda a família de motores boxer da marca: a curva de Consumo de Combustível em Carga Parcial (\textit{Part Throttle Fuel Consumption}) e a curva de Desempenho ao Nível do Mar e em Altitude (\textit{Sea Level and Altitude Performance}). A análise a seguir toma como referência metodológica as Figuras 3-5 e 3-21 do manual oficial do IO-360 \cite{io360manual}, o motor de 200 HP que, embora não seja diretamente empregado em aeronaves agrícolas da faixa de interesse, compartilha a mesma família construtiva e curvas funcionalmente equivalentes (escalonadas para o número de cilindros e cilindrada) com os motores O-540/IO-540 e IO-720, dominantes na faixa de 1 a 2{,}5 t. \autoeval{Essa metodologia foi adotada por restrição de acesso público à edição equivalente do manual da série O-540 (curvas detalhadas de Part Throttle e altitude/performance individuais por sufixo de motor).}

\subsubsection{Curva de Consumo de Combustível em Carga Parcial}

A curva de consumo de combustível em carga parcial relaciona, para diferentes regimes de RPM (1800 a 2700 RPM), o consumo em libras por hora (lb/h) com a potência efetiva entregue ao freio (\textit{Actual Brake Horsepower}). 

\textbf{COLOCAR OS GRAFICOS AQUI}

O gráfico apresenta duas famílias de curvas:

\begin{itemize}
    \item \textbf{Mistura de Melhor Potência (\textit{Best Power}):} Razão ar/combustível ligeiramente rica (em torno de 12{,}5:1), maximizando o torque entregue. Adequada às fases de decolagem, ganho de altitude e manobras de retorno em pulverização agrícola.
    \item \textbf{Mistura de Melhor Economia (\textit{Best Economy}):} Razão ar/combustível próxima à estequiométrica (cerca de 14{,}7:1) ou ligeiramente pobre. Adequada ao cruzeiro e a deslocamentos, reduzindo o consumo específico em até 15\%.
\end{itemize}

Para o IO-360 (200 HP), os dados extraídos da curva indicam consumo de aproximadamente \SI{85}{lb/h} a \SI{170}{HP} (75\% da potência nominal, regime de cruzeiro típico) em mistura \textit{best power}, e \SI{75}{lb/h} no mesmo regime em mistura \textit{best economy}. Esses valores correspondem a um BSFC de aproximadamente \SI{0.50}{lb/(hp \cdot h)} e \SI{0.44}{lb/(hp \cdot h)} respectivamente, equivalentes a \SI{304}{g/(kW \cdot h)} e \SI{268}{g/(kW \cdot h)} \cite{io360manual}.

\textbf{Extrapolação para a Família O-540 / IO-540 / IO-720:} A topologia da curva é estruturalmente idêntica para os demais motores da família, com o eixo de potência reescalado. Para o IO-540-K1J5 (320 HP do Ipanema EMB-202A), o consumo em cruzeiro a 75\% (\SI{240}{HP}) situa-se em \SI{130}{lb/h} em \textit{best power} e \SI{115}{lb/h} em \textit{best economy}; para o IO-720-D1CD (400 HP do Pawnee Brave 375), os valores escalam a aproximadamente \SI{170}{lb/h} e \SI{150}{lb/h} respectivamente \cite{lycomingmanual, lycoming_io720}.

\subsubsection{Curva de Desempenho ao Nível do Mar e em Altitude (Figura 3-21 do Manual)}

A segunda curva crítica relaciona a potência disponível com a altitude pressão (em milhares de pés) e a pressão de admissão (\textit{Manifold Pressure}, em polegadas de mercúrio). 

\textbf{COLOCAR OS GRAFICOS AQUI}

O gráfico é dividido em duas regiões:

\begin{itemize}
    \item \textbf{Quadrante Inferior --- Desempenho ao Nível do Mar:} Mostra a potência ao freio em função da pressão de admissão (variável de \SI{18}{inHg} a \SI{29}{inHg}) e do RPM. Define o envelope operacional em decolagem.
    \item \textbf{Quadrante Superior --- Desempenho em Altitude:} Mostra a perda de potência à medida que a altitude pressão aumenta (de 0 a \SI{24000}{ft}). Para um motor naturalmente aspirado (sem turbocompressor), a perda é aproximadamente linear com a densidade do ar, atingindo cerca de 75\% da potência nominal a \SI{8000}{ft} (\SI{2440}{m}) e 60\% a \SI{14000}{ft} (\SI{4270}{m}) \cite{io360manual}.
\end{itemize}

\textbf{Implicação para operações brasileiras:} Pistas agrícolas localizadas em planaltos (Goiás, Distrito Federal, sul de Minas Gerais, oeste da Bahia) frequentemente situam-se entre \SI{800}{m} e \SI{1200}{m} de elevação, com temperaturas operacionais entre \SI{30}{\degreeCelsius} e \SI{38}{\degreeCelsius} no auge da safra. Aplicando a regra prática de que motores aeronáuticos naturalmente aspirados perdem aproximadamente 3{,}5\% da potência por \SI{1000}{ft} de altitude-densidade \cite{cfisteph_density}, e considerando que dias quentes em pistas de planalto podem elevar a altitude-densidade para a faixa de \SI{4000}{ft} a \SI{8000}{ft} (cerca de \SI{1200}{m} a \SI{2400}{m}), a redução resultante de potência situa-se na faixa de aproximadamente 14\% a 28\%. Essa perda justifica a especificação de margens de potência instalada generosas sobre o MTOW e explica a opção do Cessna AgHusky pelo motor turbocompressor TSIO-520-T, a única variante que utiliza turbocompressão na faixa estudada \cite{cessna_188}.

\subsection{Especificações Comparativas}

A Tabela \ref{tab:motores_faixa} apresenta o perfil técnico consolidado dos motores das familias certificadas aplicáveis às plataformas de referência identificadas na Tabela \ref{tab:plataformas_referencia}.

\begin{table}[H]
\centering
\caption{Perfil Comparativo de Motores Aeronáuticos de Pistão para Aeronaves Agrícolas de 1 a 2{,}5 toneladas}
\label{tab:motores_faixa}
\footnotesize
\renewcommand{\arraystretch}{1.3}
\begin{tabularx}{\textwidth}{@{} 
    >{\hsize=1.0\hsize\raggedright\arraybackslash}X 
    c 
    >{\hsize=0.8\hsize\raggedright\arraybackslash}X 
    c 
    c 
    >{\hsize=1.2\hsize\raggedright\arraybackslash}X 
    @{}}
\toprule
\textbf{Modelo} & \textbf{Cil. / Desloc.} & \textbf{Potência} & \textbf{Peso Seco} & \textbf{TBO} & \textbf{Aplicação Típica (MTOW)} \\ \midrule
Lycoming O-540-B2B5 & 6 / 8{,}9 L & \SI{235}{HP} @ 2575 RPM & $\sim$\SI{179}{kg} & \SI{2000}{h} & PA-25 Pawnee B (1{,}3 t) \\
\addlinespace
Lycoming O-540-E4B5 & 6 / 8{,}9 L & \SI{260}{HP} @ 2700 RPM & $\sim$\SI{181}{kg} & \SI{2000}{h} & PA-25 Pawnee C/D (1{,}3 t) \\
\addlinespace
Continental IO-520-D & 6 / 8{,}5 L & \SI{300}{HP} @ 2850 RPM & $\sim$\SI{208}{kg} & \SI{1700}{h} & Cessna 188 AgWagon/AgTruck (1{,}5--2{,}0 t) \\
\addlinespace
Continental TSIO-520-T & 6 / 8{,}5 L (turbo) & \SI{310}{HP} @ 2700 RPM & $\sim$\SI{216}{kg} & \SI{1400}{h} & Cessna AgHusky (2{,}0 t) \\
\addlinespace
Lycoming IO-540-K1D5 / K1G5 / K1J5 & 6 / 8{,}9 L & \SI{300}{HP}/\SI{320}{HP} @ 2700 RPM & $\sim$\SI{200}{kg} & \SI{2000}{h} & Ipanema 201/202 (1{,}5--1{,}8 t), Brave 300 (2{,}0 t) \\
\addlinespace
Lycoming IO-720-D1CD & 8 / 11{,}8 L & \SI{375}{HP}--\SI{400}{HP} @ 2650 RPM & $\sim$\SI{253}{kg} & 1.800--\SI{2000}{h} & PA-36 Brave 375/400 (2{,}2 t) \\ \bottomrule
\multicolumn{6}{l}{\footnotesize Compilação do autor a partir de \cite{lycomingmanual, io360manual, lycoming_io720, continental_io520, embraer202, piper_pa25, cessna_188, piper_pa36}.}
\end{tabularx}
\end{table}

Da Tabela \ref{tab:motores_faixa} é possível inferir três observações relevantes para o projeto AGRO-VANT:

\begin{enumerate}
    \item \textbf{Densidade de potência:} A razão potência/peso seco varia de aproximadamente 1{,}3 HP/kg (O-540-B2B5) a 1{,}6 HP/kg (IO-540-K1J5 e IO-720), com pouca variação dentro da faixa.
    \item \textbf{TBO competitivo:} Todos os motores apresentam TBO igual ou superior a \SI{1700}{h}, com a maior parte da família Lycoming O-540 atingindo \SI{2000}{h} \cite{lycomingmanual}. \autoeval{Cálculo do autor: a uma utilização média anual estimada em 300 h/ano (consistente com a janela operacional sazonal típica de safra brasileira), esse TBO equivale a aproximadamente 6{,}5 anos de operação entre revisões gerais.}
    \item \textbf{Convergência no IO-540:} A família Lycoming IO-540 cobre praticamente toda a faixa central (1{,}3--2{,}0 t) com variantes do mesmo motor base, oferecendo cadeia logística unificada e disponibilidade de peças. Esse fato sustenta a recomendação do IO-540 como motor de referência para o AGRO-VANT em sua configuração-alvo.
\end{enumerate}

\section{Avaliação de Vetores Alternativos: Tração Elétrica e Hibridização}

A demanda por sistemas de menor impacto ambiental impulsionou pesquisas em plataformas puramente elétricas e híbridas. Esta seção avalia a aplicabilidade dessas tecnologias à faixa de 1 a 2{,}5 toneladas, considerando o balanço energético e as restrições de carga útil agroquímica.

\subsection{Tração Puramente Elétrica}

A inviabilidade da tração puramente elétrica para aeronaves agrícolas de 1 a 2{,}5 t é demonstrada quantitativamente pelo balanço de massa de baterias. Tomando como caso de referência o IO-540-K1J5 do EMB-202 Ipanema (\SI{320}{HP} $\approx$ \SI{239}{kW}) operando por 2 horas em regime de cruzeiro a 75\% da potência nominal (\SI{179}{kW}), e adotando como densidade de energia da bateria os valores limítrofes da literatura: Li-Ion atual em torno de \SI{250}{Wh/kg} (próximo ao limite superior dos sistemas embarcados comercialmente disponíveis \cite{battery_aviation_review, aerospace_america_batteries}) e Li-S em desenvolvimento em torno de \SI{500}{Wh/kg} (limite teórico próximo da fronteira ainda não realizada comercialmente \cite{battery_aviation_review}), o cálculo conduzido pelo autor é o seguinte:

\begin{itemize}
    \item \textbf{Energia necessária:} $E = 179\,\si{kW} \times 2\,\si{h} = 358\,\si{kWh}$
    \item \textbf{Massa de baterias Li-Ion (\SI{250}{Wh/kg}):} $m_{bat} = 358\,\si{kWh} / 0{,}25\,\si{kWh/kg} = 1432\,\si{kg}$
    \item \textbf{Massa com tecnologia emergente Li-S (\SI{500}{Wh/kg}):} $m_{bat} = 716\,\si{kg}$
\end{itemize}

Com MTOW de \SI{1800}{kg} e peso vazio aproximado de \SI{1100}{kg}, a margem disponível para carga útil + combustível + baterias é de cerca de \SI{700}{kg}. Mesmo no cenário otimista com baterias Li-S, a massa de baterias consumiria toda a margem operacional, eliminando totalmente a carga útil agroquímica. Para o Piper Brave 375 (\SI{400}{HP}, MTOW \SI{2177}{kg}), o cálculo resulta em valores ainda mais desfavoráveis.

\subsection{Trações Híbridas}

Diferentemente do vetor puramente elétrico, as arquiteturas híbridas merecem avaliação mais detalhada, dado que poderiam, em tese, reduzir o consumo de combustível e ampliar a redundância do sistema propulsivo.

\subsubsection{Modelo Híbrido em Série}
Nessa configuração, o motor a combustão atua exclusivamente como gerador, enquanto motores elétricos realizam a tração da hélice.

\begin{itemize}
    \item \textbf{Funcionamento:} O propulsor de combustão fornece energia para os Controladores Eletrônicos de Velocidade (ESC) ou para a recarga de um banco de baterias intermediário, operando em regime estacionário de RPM otimizado.
    \item \textbf{Aplicação e Desvantagens:} Essa topologia aumenta o peso da aeronave ao exigir o motor a combustão, o gerador e os motores elétricos de tração, resultando em um aumento de massa total. A literatura especializada indica ganhos de fuel burn na ordem de 12\% a 25\% em testes representativos, com o Diamond DA36 E-Star (configuração híbrida série) reportando até 25\% de redução em testes iniciais e 30\% em retrofits otimizados de aeronaves regionais \cite{wiki_hybrid_electric, sciencedirect_hybrid}. \autoeval{Para aeronaves agrícolas, esses ganhos marginais em economia de combustível tendem a ser superados pela perda de carga útil agroquímica decorrente do peso adicional do sistema híbrido, justificando o ceticismo do autor quanto à atratividade do modelo série para o segmento.}
\end{itemize}

\subsubsection{Modelo Híbrido Paralelo}
A arquitetura paralela acopla mecanicamente um motor a combustão e um motor elétrico ao mesmo eixo de tração da hélice.

\begin{itemize}
    \item \textbf{Auxílio Trativo e Regeneração:} O motor elétrico auxilia em momentos críticos (decolagem totalmente carregado, ganho de altitude) e pode atuar como gerador em regime de cruzeiro, recarregando o banco de baterias com o excesso de potência mecânica do ICE.
    \item \textbf{Vantagens Potenciais:} A integração permite o uso de um motor a combustão menor (e mais leve) do que aquele que seria necessário para a decolagem isolada. \autoeval{Como exemplo conceitual proposto pelo autor: para uma aeronave de 1{,}8 t, seria viável substituir um IO-540 de \SI{320}{HP} por um motor de \SI{200}{HP} complementado por um motor elétrico de \SI{120}{HP} dedicado às fases críticas (decolagem e subida).} A literatura aponta ganhos típicos de fuel burn entre 12\% e 30\% em arquiteturas paralelas, dependendo do grau de hibridização e do perfil de missão \cite{sciencedirect_hybrid, fuel_economy_hybrid}; em demonstrações comerciais, a Ampaire reportou economias de até 40\% no custo de combustível em rotas de curta distância \cite{ampaire_demo}.
    \item \textbf{Limitações:} Ainda em fase experimental para aeronaves certificadas. Programas como o magniX e o Ampaire EEL exploram retrofits híbridos em aeronaves leves, com demonstrações operacionais em curso desde 2020 \cite{ampaire_demo, wiki_hybrid_electric}, mas a certificação sob FAR 23 / RBAC 23 ainda não foi formalizada para uso agrícola. Os custos atuais de integração (motor elétrico aeronáutico certificado + ESC + banco de baterias certificado + sistema de gerenciamento) são proibitivos para o segmento agrícola, que opera com margens estreitas.
\end{itemize}

\section{Desempenho de Hélices e Seleção de Materiais}

A eficiência do sistema propulsivo depende da seleção adequada das hélices para a transmissão do torque. Em aeronaves agrícolas certificadas da faixa de 1 a 2{,}5 t, a hélice é um componente certificado individualmente (sob FAR 35 / RBAC 35), tipicamente fabricado por especialistas como Hartzell, McCauley ou MT-Propeller, e o seu projeto considera carregamentos aeroelásticos complexos. 

\subsection{Hélices de Velocidade Constante (Passo Variável)}
Em aeronaves da faixa de interesse, o uso de hélices de velocidade constante é padrão para otimizar o ponto de operação do motor em regimes distintos (decolagem, cruzeiro, pulverização). O Ipanema EMB-202 utiliza uma Hartzell tripá de velocidade constante \cite{embraer202, embraer202_wiki}; o PA-25 Pawnee em suas variantes mais recentes utiliza hélice McCauley ou Hartzell de velocidade constante \cite{piper_pa25}. O mecanismo de variação de passo é hidráulico, atuado por um governador (\textit{constant-speed unit}) que mantém o RPM-alvo independentemente da pressão de admissão.

A operação a velocidade constante permite explorar com eficácia as curvas de Part Throttle Fuel Consumption analisadas na Seção \ref{sec:curvas_lycoming}: durante o cruzeiro, o piloto reduz o RPM e a pressão de admissão, operando o motor próximo da curva de \textit{best economy}, enquanto na decolagem o governador permite que o motor atinja o RPM máximo na curva de \textit{best power} \cite{lycomingmanual, io360manual}. \autoeval{Valores típicos extraídos das curvas e da prática operacional são da ordem de 2.300 RPM e 22 inHg em cruzeiro econômico, contra 2.700 RPM em decolagem para os motores da família IO-540; a confirmação dos valores exatos para um motor específico exige consulta ao Pilot's Operating Handbook (POH) da aeronave correspondente.}

\subsection{Engenharia de Materiais e Seleção da Hélice}
A escolha do material da hélice afeta a eficiência aerodinâmica, a integridade estrutural, a propagação de vibrações ao \textit{frame} e a tolerância a danos por impacto com material vegetal (frequente em voos rasantes de pulverização). Em aeronaves certificadas, a escolha é determinada pelo fabricante da hélice em consonância com o motor.

Para motores aeronáuticos de pistão com torque pulsante (todos os boxer 4T da faixa de interesse), hélices de \textbf{alumínio forjado} (variantes Hartzell HC-C2YK, HC-C3YR e similares) são a opção tradicional, oferecendo equilíbrio entre rigidez torcional, capacidade de absorção de vibrações de baixa frequência e tolerância a deformação plástica em caso de impacto leve com vegetação. Hélices de \textbf{compósito} (CFRP encapsulado, com matriz de ataque metálico, como exemplo: Hartzell ASC-II) oferecem redução de massa relevante em relação ao alumínio --- o fabricante reporta valores da ordem de 20\% em linhas como a Hartzell Trailblazer \cite{hartzell_trailblazer, hartzell_composite_faq} ---, mas exigem balanceamento dinâmico cuidadoso para evitar transmissão de vibrações de alta frequência aos rolamentos do virabrequim. Hélices de \textbf{madeira laminada} (com laminação de Sitka spruce ou bétula reforçadas com revestimento de fibra de vidro) foram historicamente utilizadas em aeronaves agrícolas pioneiras (Pawnee 150), mas hoje têm uso restrito a aviões de aviação geral leves.

\begin{table}[H]
\centering
\caption{Comparativo de Materiais de Hélices Aplicáveis à Faixa de 1--2{,}5 t}
\label{tab:helices}
\footnotesize
\renewcommand{\arraystretch}{1.3}
\begin{tabularx}{\textwidth}{@{} 
    >{\hsize=1.1\hsize\raggedright\arraybackslash}X 
    c 
    c 
    >{\hsize=0.7\hsize\raggedright\arraybackslash}X 
    >{\hsize=0.9\hsize\raggedright\arraybackslash}X 
    >{\hsize=1.3\hsize\raggedright\arraybackslash}X 
    @{}}
\toprule
\textbf{Material} & \textbf{Rigidez} & \textbf{Peso} & \textbf{Vibração} & \textbf{Tolerância a Impacto} & \textbf{Aplicação na Faixa} \\ \midrule
Alumínio Forjado & Alta & Alto & Boa absorção & Muito Alta (deformação plástica) & Ipanema, Pawnee, AgWagon, Brave 300 \\
\addlinespace
Compósito CFRP encapsulado & Muito Alta & Médio & Transmite alta freq. & Alta (delaminação localizada) & Brave 375 \\
\addlinespace
Madeira Laminada & Média & Médio & Excelente absorção & Baixa (quebra segura) & Pawnee 150 (histórico, fora de produção) \\ \bottomrule
\multicolumn{6}{l}{\footnotesize Compilação do autor com referência a \cite{tiruvenkadam2024} e catálogos Hartzell, McCauley.}
\end{tabularx}
\end{table}

\section{Sensoriamento, Telemetria e Monitoramento}

Para garantir a segurança operacional e o monitoramento detalhado de aeronaves agrícolas tripuladas, a arquitetura eletrônica embarcada cobre quatro funções principais: (i) instrumentação do motor (CHT, EGT, fluxo de combustível, pressão de óleo, RPM); (ii) navegação e geo-referenciamento de pulverização; (iii) controle automático da vazão de aplicação; e (iv) comunicação rádio com a estação de solo e demais aeronaves do espaço aéreo.

\subsection{Protocolos de Comunicação e Controle}
\begin{itemize}
    \item \textbf{Barramento Interno (CAN Bus / ARINC 825):} Em motores certificados modernos (Lycoming iE2), o gerenciamento eletrônico utiliza barramento CAN Bus em variante aeronáutica (padrão ARINC 825), interligando FADEC, sensores de motor e o painel EICAS (\textit{Engine Indication and Crew Alerting System}) \cite{lycomingmanual, lycoming_io720}.
    
    \item \textbf{Sistemas Específicos de Aviação Agrícola:} Plataformas agrícolas modernas utilizam sistemas GPS dedicados de provedores especializados como AG-NAV \cite{agnav_official}, Satloc \cite{satloc_intelliflow} e TracMap \cite{tracmap_aviation} para geo-referenciamento da deriva de aplicação e controle automático de vazão (\textit{flow control}) acoplado à velocidade no solo (GS) lida pelo GPS. Esses sistemas permitem aplicação variável de defensivos por zona da lavoura (\textit{variable-rate application}) baseada em mapas de prescrição (PMAPs) \cite{satloc_intelliflow, kirk_1996_satloc}, reduzindo desperdício e impacto ambiental ao corrigir variações na vazão durante passes upwind/downwind em decorrência de mudanças na velocidade no solo \cite{kirk_1996_satloc}.
    
    \item \textbf{Aviônica de Cabine:} Os dados anteriormente enviado para a cabine por meio Sistemas de Voo Eletrônicos (EFIS) como Garmin G600 TXi ou Avidyne IFD-540, fornecendo informações de motor, navegação e geo-referenciamento de pulverização podem ser enviados Estação de Controle em Solo (GCS - Ground Control Station). Porém alguns exemplos de aeronaves clássicas (Pawnee, AgWagon, Ipanema EMB-201) ainda utilizam instrumentação analógica, dificultando a integração com sistemas de comunicação modernos para drones não tripulados.
\end{itemize}

\subsection{Monitoramento de Fluidos e Temperatura}
A precisão na autonomia e a proteção termomecânica demandam sensores de leitura direta:

\begin{itemize}
    \item \textbf{Sensor de Fluxo (Flowmeter):} Mede o consumo volumétrico instantâneo na linha de combustível, apresentando maior precisão que boias de nível mecânicas, as quais sofrem interferência por inércia durante o voo. Em motores Lycoming IO-540, o sensor de fluxo é parte integrante do indicador de mistura (\textit{Fuel Flow / EGT analyzer}) acoplado a sistemas certificados como o JPI EDM-900/930, que substitui a instrumentação analógica de motor e integra RPM, MAP, EGT/CHT por cilindro, fluxo e pressão de combustível e óleo em uma única tela \cite{jpi_edm900, jpi_edm930_manual}.
    
    \item \textbf{Termopares (Tipo K) por cilindro:} São essenciais para o gerenciamento da saúde do motor a combustão. Em motores de pistão das famílias O-540/IO-540/IO-720, a instrumentação moderna inclui CHT (\textit{Cylinder Head Temperature}) e EGT (\textit{Exhaust Gas Temperature}) individuais para cada cilindro (6 ou 8 sensores), permitindo diagnóstico precoce de cilindros operando fora da estequiometria ideal:
    \begin{itemize}
        \item \textit{CHT:} Limite máximo de \SI{500}{\degree F} (\SI{260}{\degreeCelsius}) no IO-540 \cite{lycomingmanual}; valores acima indicam refrigeração comprometida ou \textit{detonation}.
        \item \textit{EGT:} Variações anômalas indicam mistura ar/combustível fora da estequiometria ideal, permitindo ajustes manuais ou correções automatizadas via FADEC. O monitoramento de EGT por cilindro é a base para a técnica de operação \textit{Lean of Peak} (LOP), que segundo a literatura técnica de aviação geral pode reduzir o consumo de combustível em torno de 20\% em cruzeiro, ao custo de uma redução modesta (cerca de 5\%) na velocidade \cite{aopa_lop}. Em campo, operadores reportam tipicamente economias da ordem de 2 a 3 galões/hora em motores grandes (Continental e Lycoming) operando LOP com injetores balanceados \cite{aviation_consumer_gami}.
        \item Em alguns casos a adaptação do ponto de ignição do combustivel pela vela quando em velocidades de cruzeiro tornaria o regime econômico mais eficiente, mas isso exigiria modificações no sistema de ignição.
    \end{itemize}
\end{itemize}

\subsection{Digital Twins}

A integração de \textit{Digital Twins} surge como tecnologia habilitadora para monitoramento em tempo real, manutenção preditiva e simulação de condições de contorno realistas. Para o AGRO-VANT, propõe-se a estruturação de um \textit{digital twin} composto por três camadas: (i) um modelo termodinâmico do conjunto motopropulsor calibrado por dados de bancada e pelas curvas de \textit{Part Throttle Fuel Consumption} apresentadas na Seção \ref{sec:curvas_lycoming}; (ii) um modelo multi-corpo da estrutura da aeronave; e (iii) um modelo de degradação de componentes (fadiga acumulada na hélice, desgaste de cilindros) alimentado pela telemetria embarcada. Essa abordagem está alinhada com práticas emergentes da Embraer e dos fabricantes de motores aeronáuticos para frotas em operação, e pode subsidiar políticas de manutenção baseada em condição (CBM --- \textit{Condition-Based Maintenance}). \autoeval{A quantificação de ganhos específicos (e.g., extensão do TBO efetivo) depende de validação experimental ainda não conduzida no escopo deste estudo, configurando linha de pesquisa para os próximos estágios do projeto.}

\section{Referências Normativas e Requisitos de Certificação}

A operação e certificação de aeronaves agrícolas tripuladas na faixa de 1 a 2{,}5 toneladas no Brasil estão sob jurisdição da Agência Nacional de Aviação Civil (ANAC) e regidas por um conjunto de regulamentos específicos:

\begin{itemize}
    \item \textbf{RBAC 23:} Estabelece os requisitos de aeronavegabilidade para aeronaves de categoria normal (MTOW $\leq$ \SI{8618}{kg}) \cite{anacrbac23}. É o regulamento aplicável a todas as aeronaves de referência analisadas neste trabalho (Pawnee, AgWagon, Ipanema, Pawnee Brave). Equivalente ao FAR 23 (FAA) e CS-23 (EASA).
    
    \item \textbf{RBAC 33:} Regulamento de certificação de motores aeronáuticos a pistão. Todos os motores Lycoming O-540, IO-540 e IO-720, bem como os Continental IO-520 e TSIO-520 mencionados neste trabalho, possuem certificações da FAA sob o predecessor CAR 13 ou RBAC 33, com type certificates específicos (e.g., TC 1E4 para a família IO-540, TC E-295 para a O-540).
    
    \item \textbf{RBAC 35:} Regulamento de certificação de hélices. As hélices Hartzell e McCauley utilizadas nas aeronaves de referência possuem type certificates próprios (e.g., Hartzell HC-C2YK sob TC P-924).
    
    \item \textbf{RBAC 137 --- Emenda nº 05 (Resolução nº 716/2023):} Em vigor desde 2 de outubro de 2023, estabelece o cadastro e os requisitos operacionais para operações aeroagrícolas no Brasil \cite{anac137}. A principal mudança da Emenda 05 foi a substituição do antigo Certificado de Operador Aéreo (COA) pelo Cadastro de Operador Aeroagrícola (CDAG), simplificando o processo de habilitação. A frota brasileira ao abrigo do regulamento conta com aproximadamente \num{2500} aeronaves, em sua maioria Ipanemas e variantes \cite{anac137}.
    
    \item \textbf{IS 137.201-001E:} Instrução Suplementar que regulamenta o uso do etanol em aeronaves agrícolas, exigindo Autorização Especial de Voo (AEV) emitida sob a Subparte 21.175 do RBAC 21 \cite{anacis137}. Aplicável especificamente ao caso do EMB-202A.
    
    \item \textbf{Normas ABNT aplicáveis:} A ABNT NBR ISO 9001 é referência para sistemas de qualidade em fabricantes; a NBR 14785 trata da segurança operacional na aviação agrícola; e a NBR ISO 21384 (em adaptação a partir do ISO/TC 20/SC 16) cobre futuras aeronaves não tripuladas. Para o segmento tripulado tradicional da faixa, a referência normativa principal permanece sendo os RBACs.
    
    \item \textbf{FAR Part 137 (FAA) e CS-VLA / CS-23 (EASA):} Regulamentos internacionais análogos, relevantes para certificação cruzada via acordos bilaterais (e.g., o BASA Brasil-EUA, que reconhece automaticamente type certificates da FAA na frota brasileira).
\end{itemize}

A escolha do grupo motopropulsor é diretamente condicionada por esses regulamentos. Motores com FAA Type Certificate (toda a família O-540, IO-540, IO-720, IO-520, TSIO-520 listada na Tabela \ref{tab:motores_faixa}) atendem automaticamente aos requisitos de aeronavegabilidade brasileiros via acordos bilaterais, com homologação simplificada. Motores não certificados (motores automotivos convertidos, motores 2T de uso industrial, motores BLDC de aplicação \textit{drone}) somente podem ser empregados sob regime de aeronave experimental (Subparte 21.191 do RBAC 21) com restrições operacionais e proibições de transporte oneroso.

\section{Conclusão e Combinações Recomendadas}

O estudo demonstra que a faixa de aeronaves agrícolas tripuladas de 1 a 2{,}5 toneladas apresenta convergência tecnológica notável em motores aeronáuticos de pistão certificados, de configuração boxer horizontalmente oposta, com potências entre \SI{235}{HP} e \SI{400}{HP}. A análise das curvas oficiais Lycoming (Seção \ref{sec:curvas_lycoming}), o levantamento das plataformas de referência (Tabela \ref{tab:plataformas_referencia}) e a comparação direta dos motores aplicáveis (Tabela \ref{tab:motores_faixa}) convergem para um conjunto de recomendações estruturadas por cenário de massa máxima de decolagem.

\subsection{Cenários Operacionais Recomendados}

\begin{table}[H]
\centering
\caption{Recomendações de Combinações Propulsivas por Cenário de Missão na Faixa 1--2{,}5 t}
\label{tab:recomendacoes}
\footnotesize
\renewcommand{\arraystretch}{1.3}
\begin{tabularx}{\textwidth}{@{} 
    >{\hsize=0.9\hsize\raggedright\arraybackslash}X 
    >{\hsize=1.0\hsize\raggedright\arraybackslash}X 
    >{\hsize=0.9\hsize\raggedright\arraybackslash}X 
    >{\hsize=1.2\hsize\raggedright\arraybackslash}X 
    @{}}
\toprule
\textbf{Cenário (MTOW)} & \textbf{Motor + Hélice} & \textbf{Aviônica} & \textbf{Combustível / Observação} \\ \midrule
\textbf{Cenário A:}\newline AGRO-VANT 1{,}0--1{,}5 t,\newline missão $\leq$ \SI{2}{h},\newline pulverização leve & Lycoming O-540-E4B5 (260 HP) ou IO-540-K1D5 (300 HP) + Hartzell HC-C2YK bipá Al forjado, vel. constante & EFIS Garmin G500 TXi + DGPS Satloc / TracMap, sensor JPI EDM-900 com CHT/EGT por cilindro & AvGas 100LL; alternativa etanol via IS 137.201-001E após retrofit \\
\addlinespace
\textbf{Cenário B (Recomendado):}\newline AGRO-VANT 1{,}5--2{,}0 t,\newline missão $\leq$ \SI{3}{h},\newline equivalente Ipanema/Brave 300 & Lycoming IO-540-K1J5 (320 HP) + Hartzell HC-C3YR tripá Al forjado, vel. constante & EFIS Garmin G600 TXi + DGPS AGNAV, EICAS dedicado, retrofit para FADEC Lycoming iE2 quando disponível & Etanol via IS 137.201-001E (paradigma EMB-202A): vantagem econômica e ambiental para o cenário brasileiro \\
\addlinespace
\textbf{Cenário C:}\newline AGRO-VANT 2{,}0--2{,}5 t,\newline alta produtividade,\newline grandes hoppers & Lycoming IO-720-D1CD (375--400 HP) + Hartzell HC-E4N5L quadripá compósita encapsulada, vel. constante & EFIS Garmin G1000 NXi + DGPS RTK, ADS-B Out, manutenção baseada em condição (CBM) & AvGas 100LL; alta densidade de potência justifica IO-720 sobre IO-540 \\
\addlinespace
\textbf{Cenário D:}\newline Hot-and-High (planaltos, dias quentes) em qualquer MTOW da faixa & Continental TSIO-520-T turbocomprimido (310 HP) com waste-gate variável + hélice tripá vel. constante & Aviônica padrão + sensor de pressão de \textit{manifold} para gerenciamento do turbocompressor & AvGas 100LL; \textit{boost} compensa perda de densidade do ar (justifica AgHusky em planaltos) \\ \bottomrule
\end{tabularx}
\end{table}

\subsection{Síntese das Recomendações}

\begin{enumerate}
    \item \textbf{Para a configuração base do AGRO-VANT (MTOW $\sim$\SI{1500}{kg}--\SI{2000}{kg}):} A recomendação primária é a adoção do conjunto Lycoming IO-540-K1J5 + Hartzell tripá HC-C3YR de velocidade constante (Cenário B), replicando o paradigma consolidado do EMB-202 Ipanema. Essa combinação oferece o melhor compromisso entre custo de aquisição, disponibilidade de assistência técnica no Brasil, suporte à operação com etanol (apoiando a agenda nacional de descarbonização) e conformidade direta com o RBAC 23 e RBAC 137. O peso seco do motor (\SI{200}{kg}) e seu envelope dimensional são compatíveis com a estrutura típica de aeronaves agrícolas leves.
    
    \item \textbf{Para configurações com MTOW superior a \SI{2000}{kg}:} A migração para o Lycoming IO-720-D1CD (400 HP, 8 cilindros) é justificada quando o aumento de massa total exige potência adicional não fornecida pelo IO-540. O acréscimo de aproximadamente \SI{53}{kg} no peso seco do motor é compensado pelo ganho de potência de 25\% e pela maior margem de potência em densidade de altitude.
    
    \item \textbf{Para operações em alta densidade de altitude:} Considerar o Continental TSIO-520-T turbocomprimido (caso AgHusky) como alternativa, em particular para aeronaves operando frequentemente em pistas acima de \SI{1000}{m} de elevação e temperaturas acima de \SI{32}{\degreeCelsius}. A turbocompressão permite manter potência nominal em altitudes onde motores naturalmente aspirados perdem 10\% a 15\% da potência.
    
    \item \textbf{Para a aviônica e gestão de propulsão:} Recomenda-se a adoção de instrumentação digital integrada (EFIS Garmin TXi ou equivalente) com monitoramento CHT/EGT por cilindro, sensor de fluxo de combustível e DGPS de uso agrícola integrado. A arquitetura de \textit{digital twin} para manutenção baseada em condição representa caminho de modernização recomendável.
    
    \item \textbf{Sobre tração elétrica e hibridização:} A tração puramente elétrica é tecnicamente inviável para a faixa de 1 a 2{,}5 t em horizonte de 10 a 15 anos, dado o balanço de massa de baterias atual e previsível. A hibridização paralela representa linha de pesquisa promissora, mas ainda dependente de avanços em certificação aeronáutica de motores elétricos e baterias.
\end{enumerate}

\newpage

\begin{thebibliography}{99}

    % Fundamentos de Propulsão Aeroespacial e Combustão
    \bibitem{elsayed2016}
    El-Sayed, A. F. (2016). \textit{Fundamentals of Aircraft and Rocket Propulsion}. London: Springer.

    \bibitem{heywood1988}
    Heywood, J. B. (1988). \textit{Internal Combustion Engine Fundamentals}. New York: McGraw-Hill Education.

    % Materiais de Hélices
    \bibitem{tiruvenkadam2024}
    Tiruvenkadam, N., et al. (2024). \textit{Investigation of Structural and Thermal Analysis of Drone Propeller Materials}. Journal of Physics: Conference Series, 2925, 012002. IOP Publishing. Disponível em: \url{https://doi.org/10.1088/1742-6596/2925/1/012002}

    % ==== Manuais e dados oficiais de fabricantes de motores aeronáuticos ====

    \bibitem{io360manual}
    Lycoming Engines. \textit{Operator's Manual --- O-360, HO-360, IO-360, AIO-360, HIO-360 \& TIO-360 Series Aircraft Engines}, 8th Edition, Part No. 60297-12, Approved by FAA. Williamsport, PA: October 2005. Figuras 3-5 (Part Throttle Fuel Consumption --- IO-360-A,-C,-D,-J,-K; AIO-360 Series) e 3-21 (Sea Level and Altitude Performance --- mesma série).

    \bibitem{lycomingmanual}
    Lycoming Engines. \textit{Operator's Manual --- O-540 and IO-540 Series Aircraft Engines}, 4th Edition, Approved by FAA. Disponível em: \url{https://www.lycoming.com/sites/default/files/file/2025-02/60297-10\%20-\%20O-540\%20and\%20IO-540\%20Series.pdf}

    \bibitem{lycoming_io720}
    Lycoming Engines. \textit{720 Series Engines --- Technical Specifications}. Williamsport, PA. Acessado em: 12 de maio de 2026. Disponível em: \url{https://www.lycoming.com/engines}

    \bibitem{continental_io520}
    Continental Motors Inc. \textit{IO-520 and TSIO-520 Series Engine Specifications and Type Certificate Data Sheet}. Acessado em: 12 de maio de 2026.

    % ==== Aeronaves agrícolas de referência (faixa 1--2,5 t) ====

    \bibitem{embraer202}
    Embraer / Indústria Aeronáutica Neiva. \textit{EMB-200, EMB-201, EMB-202 e EMB-203 Ipanema --- Especificações Técnicas}. Botucatu, Brasil. Acessado em: 12 de maio de 2026.

    \bibitem{piper_pa25}
    Piper Aircraft Corporation. \textit{PA-25 Pawnee Type Certificate Data Sheet and Pilot's Operating Handbook}. Variantes PA-25-150 (Lycoming O-320), PA-25-235 Pawnee B/C/D (Lycoming O-540-B2B5) e PA-25-260 (Lycoming O-540-E4B5). Acessado em: 12 de maio de 2026.

    \bibitem{cessna_188}
    Cessna Aircraft Company. \textit{Model 188 Series AgWagon, AgPickup, AgTruck and AgHusky --- Type Certificate and Operating Handbook}. Wichita, KS. Acessado em: 12 de maio de 2026.

    \bibitem{piper_pa36}
    Piper Aircraft Corporation / WTA Incorporated. \textit{PA-36 Pawnee Brave Series --- Type Certificate and Pilot's Operating Handbook}. Variantes Brave 300 (Lycoming IO-540-K1G5), Brave 375 e Brave 400 (Lycoming IO-720-D1CD). Acessado em: 12 de maio de 2026.

    % ==== Referências normativas ANAC ====

    \bibitem{anac137}
    Agência Nacional de Aviação Civil (ANAC). \textit{RBAC nº 137 --- Emenda 05: Cadastro e Requisitos Operacionais: Operações Aeroagrícolas}. Resolução nº 716, de 13 de junho de 2023. Vigência: 2 de outubro de 2023.

    \bibitem{anacis137}
    Agência Nacional de Aviação Civil (ANAC). \textit{Instrução Suplementar IS 137.201-001E --- Uso de Etanol em Aeronaves Agrícolas}. Aprovada pela Portaria nº 217/SAR, de 26 de janeiro de 2022.

    \bibitem{anacrbac23}
    Agência Nacional de Aviação Civil (ANAC). \textit{RBAC nº 23 --- Requisitos de Aeronavegabilidade: Aviões Categoria Normal}. Aplicável a aeronaves com MTOW $\leq$ \SI{8618}{kg}.

    % ==== Referências adicionais (revisão de citações) ====

    \bibitem{battery_aviation_review}
    Authors collective. \textit{Battery technology for sustainable aviation: a review of current trends and future prospects}. Applied Energy / ScienceDirect, 2025. Energy densities reported: jet fuel \SI{12000}{Wh/kg}; current Li-Ion up to \SI{330}{Wh/kg}; solid-state up to \SI{400}{Wh/kg}. Disponível em: \url{https://www.sciencedirect.com/science/article/pii/S0306261925010864}

    \bibitem{aerospace_america_batteries}
    Aerospace America (AIAA). \textit{The Buzz Over Batteries}. ``Jet A, the standard kerosene fuel for commercial and military aircraft, has an energy density of 12,000 watt-hours per kilogram.'' Disponível em: \url{https://aerospaceamerica.aiaa.org/features/the-buzz-over-batteries/}

    \bibitem{nas_army_fuels}
    National Academies of Sciences, Engineering, and Medicine. \textit{Powering the U.S. Army of the Future}, Chapter 3: ``Energy Sources, Conversion Devices, and Storage''. Washington, DC: The National Academies Press, 2021. doi:10.17226/26052. ``Ethanol [...] have 69 percent [...] of the energy content per unit volume of JP8.''

    \bibitem{cfisteph_density}
    CFI Steph. \textit{Density Altitude --- Effects on Engine Power}. ``A normally aspirated aircraft engine will lose approximately 3.5 percent of its horsepower for every 1,000-foot increase in density altitude.'' Acessado em maio de 2026. Disponível em: \url{https://www.cfisteph.com/densityaltitude}

    \bibitem{embraer202_wiki}
    Wikipedia / Embraer. \textit{Embraer EMB 202 Ipanema --- General Characteristics}. ``Hopper capacity: 950 litres [...] Max takeoff weight: 1,800 kg [...] Empty weight: 1,020 kg [...] Fuel capacity: 264 litres usable.'' Disponível em: \url{https://en.wikipedia.org/wiki/Embraer_EMB_202_Ipanema}

    \bibitem{wiki_hybrid_electric}
    Wikipedia. \textit{Hybrid electric aircraft}. ``The Diamond DA36 E-Star [...] reducing fuel consumption and emissions by up to 25\%.'' Disponível em: \url{https://en.wikipedia.org/wiki/Hybrid_electric_aircraft}

    \bibitem{sciencedirect_hybrid}
    Pornet, C.; Isikveren, A. T. \textit{Conceptual design of hybrid-electric transport aircraft / Hybrid electric aircraft design with optimal power management}. ScienceDirect. ``Retrofit Do-228 found block fuel reductions of approximately 30\% for serial and parallel architectures.'' Disponível em: \url{https://www.sciencedirect.com/science/article/abs/pii/S1270963824006102}

    \bibitem{fuel_economy_hybrid}
    Voskuijl, M.; et al. \textit{Fuel economy of hybrid electric flight}. Applied Energy / ScienceDirect, 2017. ``Improvement up to 12\% in fuel consumption with respect to a non-hybrid case.'' Disponível em: \url{https://www.sciencedirect.com/science/article/abs/pii/S0306261917312631}

    \bibitem{ampaire_demo}
    Aerospace America (AIAA). \textit{Hybrid Aircraft One Step Toward the Future of Aviation}. ``Ampaire demonstrated up to 40\% fuel-cost savings'' em rotas inter-ilhas no Havaí. Disponível em: \url{https://aerospaceamerica.aiaa.org/institute/hybrid-aircraft-one-step-toward-the-future-of-aviation/}

    \bibitem{hartzell_trailblazer}
    Hartzell Propeller Inc. \textit{Trailblazer Composite Propeller --- Technical Description}. ``Materials [...] 20\% lighter than aluminum blade propellers.'' Acessado em maio de 2026.

    \bibitem{hartzell_composite_faq}
    Hartzell Propeller Inc. \textit{Composite Propellers: Why Make the Switch?} / \textit{Technical FAQ}. ``Structural composite blades offer a substantial weight reduction over aluminum propeller blades.'' Disponível em: \url{https://hartzellprop.com/composite-propellers-why-make-the-switch/}

    \bibitem{agnav_official}
    AG-NAV Inc. \textit{AG-NAV GPS Precision Navigation Systems --- Ag-Flow and Ag-Gate Application Control}. Disponível em: \url{https://www.agnav.com/}

    \bibitem{satloc_intelliflow}
    Satloc / Hemisphere GNSS. \textit{IntelliFlow 2 / IntelliFlow 3 Aerial Spray Flow Control System --- Variable Rate based on Prescription Maps (PMAPs)}. Disponível em: \url{https://www.satloc.com/products/intelliflow-2/}

    \bibitem{tracmap_aviation}
    TracMap. \textit{TracMap Aviation System with Variable Rate Flow Controller --- Press Release}. AgAir Update, 2023. Disponível em: \url{https://agairupdate.com/2023/03/03/tracmap-releases-fully-integrated-variable-rate-flow-controller-for-aviation/}

    \bibitem{kirk_1996_satloc}
    Kirk, I. W. \textit{Precision GPS Flow Control for Aerial Spray Applications}. In: Precision Agriculture Proceedings, ASA-CSSA-SSSA Books, Wiley, 1996. doi:10.2134/1996.precisionagproc3.c96. Estudo conduzido em Cessna AgHusky com sistema Satloc AIRSTAR.

    \bibitem{jpi_edm900}
    J. P. Instruments. \textit{EDM-900 Primary Engine Monitor --- Certified Replacement for Analog Engine Gages}. Inclui RPM, MAP, EGT/CHT, fluxo, pressões e temperaturas. Disponível em: \url{https://www.jpinstruments.com/}

    \bibitem{jpi_edm930_manual}
    J. P. Instruments. \textit{Pilot's Guide --- Engine Data Management EDM-930 Primary TSO}. Manual operacional. Disponível em: \url{https://www.jpinstruments.com/wp-content/uploads/2012/08/PG-EDM-930-Primary-REV-K-02-22-16-pls-Repaired.pdf}

    \bibitem{aopa_lop}
    Aircraft Owners and Pilots Association (AOPA). \textit{Technique: Lean-of-peak engine operations}. AOPA Flight Training Magazine, 2019. ``Running aircraft engines lean of peak typically reduces airspeed about 5 percent in cruise while lowering fuel consumption about 20 percent.'' Disponível em: \url{https://www.aopa.org/news-and-media/all-news/2019/may/flight-training-magazine/technique-lean-of-peak}

    \bibitem{aviation_consumer_gami}
    Aviation Consumer / AVweb. \textit{GAMIjectors: Precision Fuel Injection (18 Years)}. ``Saving generally in the range of two to three GPH when operating LOP, lower CHT, fewer fouled plugs and longer cylinder and engine life.'' Disponível em: \url{https://aviationconsumer.com/industry-news/gamijectors-precision-fuel-injection/}

\end{thebibliography}

\end{document}